
\subsubsection{Acceptance Rates}

The median per-user acceptance rate was 92.0\%, exceeding the pre-registered 70\% threshold by 22 percentage points. Overall acceptance rate (calculated at suggestion level) was 89.5\%. Total suggestions: 1,875 across 75 users.

\begin{table}[htbp]
\centering
\begin{tabular}{llll}
\toprule
Metric & Value & Target & Status \\
\midrule
Median Acceptance Rate & 92.0\% & ≥70\% & Pass \\
Overall Acceptance Rate & 89.5\% & ≥70\% & Pass \\
Total Suggestions & 1,875 & 1,000-3,000 & Pass \\
Pilot Users & 75 & 50-100 & Pass \\
\bottomrule
\end{tabular}
\end{table}

For contextual comparison, GitHub Copilot achieves 65-75\% acceptance for code generation. The 92\% rate observed in this documentation pilot may reflect lower error stakes and higher time pressure in documentation tasks, though cross-domain comparison is illustrative rather than statistically rigorous due to different application contexts, user populations, and task requirements.

\subsubsection{Cross-Domain Generalization}

Acceptance rates by dataset type showed minimal variance:

\begin{table}[htbp]
\centering
\begin{tabular}{lll}
\toprule
Dataset Type & Acceptance Rate & \# Suggestions \\
\midrule
Vision & 89.7\% & 625 \\
NLP & 89.8\% & 625 \\
Tabular & 88.8\% & 625 \\
\bottomrule
\end{tabular}
\end{table}

Kruskal-Wallis H-test: p = 0.82 (not significant). The 1.0 percentage point variance across domains indicates that a single model architecture generalizes effectively without domain-specific tuning.

\subsubsection{User Engagement}

User actions distributed as follows:

\begin{table}[htbp]
\centering
\begin{tabular}{lll}
\toprule
Action & Count & Percentage \\
\midrule
Accepted as-is & 1,177 & 62.8\% \\
Modified & 502 & 26.8\% \\
Rejected & 196 & 10.5\% \\
\bottomrule
\end{tabular}
\end{table}

The 26.8\% modification rate indicates users actively refined suggestions rather than blindly accepting them. This pattern suggests users accept helpful content into editable fields, then refine phrasing or add specifics.
